\chapter{Visits of industries}

\section{AGC - Glass industry}

\section{Aperam - Stainless steel}

\section{Biogas}

\section{Holcim Visit - Cement Industry}

% ==========================================
\subsection{Global and Regional Industrial Footprint}
% ==========================================
Holcim stands as a leading global partner for sustainable construction, operating across 43 countries on 5 continents with a workforce exceeding 45,000 employees. The European market serves as an operational spearhead for structural decarbonization, managing 38 cement plants and approximately 24,000 employees. On a global scale, circular economy integration has advanced significantly, with Holcim reaching an annual recycling rate of \num{8} million tonnes (Mt) of construction and demolition waste (CDW) in 2025 across 109 circular construction hubs.

\subsubsection{Holcim Belgium Market Dynamics and Supply Logistics}
In Belgium, Holcim’s vertically integrated local value chain operates with roughly 900 employees across 30 active sites, including:
\begin{itemize}
    \item \textbf{Binder Production:} The main integrated cement plant in Obourg (Mons) and a specialized maritime cement terminal in Antwerp.
    \item \textbf{Aggregates Extraction:} 7 limestone and chalk quarries located across strategic structural basins (e.g., Milieu in Gaurain-Ramecroix, Bruyelle, Givet, Ermitage in Deux-Acren, Leffe, Perlonjour in Naast, and Trooz).
    \item \textbf{Downstream Operations:} 8 Ready-Mix Concrete (RMX) plants (including Brussels, Aarschot, Dampremy, Gembloux, Merksem, Ghlin, Overijse, and Sint-Truiden) and 10 Mortar \& Construction Chemicals facilities (operating under brands like Cantillana and Compaktuna).
\end{itemize}

The regional macro-economic distribution reveals that the primary market demand for concrete construction is concentrated in northern Belgium and the Netherlands, meaning logistical streams are heavily oriented northward. Holcim Belgium achieves a volume throughput of \num{1.4}~Mt of cement sold annually. This binder production supports the commercialization of \num{7.6}~Mt of aggregates and massive quantities of readymix concrete, generating a resilient revenue stream while committing to a target of a \num{22}\% reduction in absolute $\text{CO}_2$ emissions per net sales relative to the 2022 baseline.

% ==========================================
\subsection{Fundamentals of Concrete and Cement Chemistry}
% ==========================================

\subsubsection{System Definitions: Concrete vs. Cement}
From an engineering perspective, it is critical to separate the low-energy downstream blending of concrete from the high-energy upstream manufacturing of cement:
\begin{enumerate}
    \item \textbf{Ready-Mix Concrete:} A 5-component heterogeneous structural system comprising binder (cement and supplementary cementitious materials), water, aggregates (graded sand and gravel/stone coarse matrix), chemical admixtures, and entrained/entrapped air pores. The water-to-binder ($w/b$) ratio dictates mechanical kinetics via the hydration chemical reaction. Downstream production is a macro-mechanical mixing operation.
    \item \textbf{Cement:} A finely ground hydraulic mineral powder acting as the continuous structural adhesive. It consists primarily of Portland clinker blended with secondary components like blast-furnace slag (GGBS), fly ash, or pulverized limestone. Manufacturing cement requires sophisticated thermo-chemical manufacturing plants capable of synthesizing artificial mineral phases.
\end{enumerate}

\subsubsection{Clinker Pyroporcessing and Thermochemistry}
Portland clinker synthesis relies on raw materials rich in calcium carbonate ($\text{CaCO}_3$, such as chalk and limestone) combined with argillaceous components supplying silica ($\text{SiO}_2$), alumina ($\text{Al}_2\text{O}_3$), and iron oxide ($\text{Fe}_2\text{O}_3$). 

The baseline process begins with quarry extraction, heavy crushing, and high-precision raw meal blending (homogenization). The material is then subjected to pyropressing. The core thermochemical reaction during the process is the calcination of calcium carbonate followed by sintering (clinkering) at material temperatures reaching up to \num{1400}~\si{\celsius} to \num{1450}~\si{\celsius} inside an industrial furnace.

An exemplary simplified chemical equation representing the formation of Alite (Tricalcium Silicate, $\text{Ca}_3\text{SiO}_5$), the primary phase responsible for early and ultimate mechanical strength in concrete, can be modeled as follows:
\begin{equation}
    3\text{CaCO}_3 + \text{SiO}_2 \xrightarrow{\Delta, \sim 1400^\circ\text{C}} \text{Ca}_3\text{SiO}_5 + 3\text{CO}_2\uparrow
\end{equation}

In cement chemistry notation (where $\text{C} = \text{CaO}$, $\text{S} = \text{SiO}_2$), this corresponds to the reaction of lime with quartz to yield $\text{C}_3\text{S}$. This reaction shows that \textbf{process-inherent $\text{CO}_2$ release is chemically unavoidable}, as it is a direct consequence of the thermal decomposition of limestone ($\text{CaCO}_3 \rightarrow \text{CaO} + \text{CO}_2$). 

Once synthesized, the incandescent clinker nodules are rapidly cooled to stabilize the hydraulic mineral phases ($\text{C}_3\text{S}$, $\text{C}_2\text{S}$, $\text{C}_3\text{A}$, $\text{C}_4\text{AF}$). Finally, clinker is interground with calcium sulfate (gypsum, $\text{CaSO}_4\cdot2\text{H}_2\text{O}$), iron sulfate, and supplementary limestone to form the final commercial cement powder.

% ==========================================
\subsection{Greenhouse Gas Emissions Profile and Decarbonization Pillars}
% ==========================================

\subsubsection{Greenhouse Gas Protocol Categorization}
As an energy-intensive and hard-to-abate industry, cement production exhibits an emissions footprint deeply rooted in its direct chemical processing operations. Approximately \num{60}\% of all cement-related carbon emissions worldwide originate strictly from the raw material process. The breakdown at Holcim satisfies standard Greenhouse Gas (GHG) protocol definitions:

\begin{table}[!ht]
\centering
\small
\begin{tabular}{llp{9cm}}
\toprule
\textbf{Emissions Category} & \textbf{Operational Allocation} & \textbf{Specific Origin / Industrial Driver} \\
\midrule
\textbf{Scope 1 (Direct)} & \textbf{52.6\%} (Total Segment) & \\
\quad -- Process Calcination & 35.5\% of total emissions & Endothermic dissociation of $\text{CaCO}_3 \rightarrow \text{CaO} + \text{CO}_2$ in the kiln.\\
\quad -- Thermal Combustion & 17.1\% of total emissions & High-temperature fuel burning required to maintain raw meal sintering temperatures.\\
\midrule
\textbf{Scope 2 (Indirect)} & \textbf{Marginal} & Electricity consumed by high-power grinding mills, raw meal fans, and ancillary quarry machinery.\\
\midrule
\textbf{Scope 3 (Value Chain)} & \textbf{Significant Balance} & Capital investments, upstream material procurement, third-party logistics, and product distribution.\\
\bottomrule
\end{tabular}
\caption{Holcim operational carbon emission breakdown structure.}
\label{tab:scopes}
\end{table}

On a global scale, built environment assets are responsible for \num{37}\% of total $\text{CO}_2$ emissions. This footprint is divided into \textbf{embodied carbon} (\num{10}\%, originating from materials procurement and the active construction phase) and \textbf{operational carbon} (\num{27}\%, driven by lifetime building heating, cooling, and power performance). Holcim is leveraging its structural materials portfolio to optimize both parameters.

\subsubsection{Strategic Decarbonization Framework: The 4 Pillars}
To transition toward a verified Net-Zero operational state by 2030, Holcim has implemented an asset transformation program structured around four strategic pillars:
\begin{enumerate}
    \item \textbf{Green Operations (Decarbonize Processes):} Transitioning industrial assets away from fossil fuels, deploying thermal energy efficiency systems, and modernizing kiln architectures.
    \item \textbf{Building Better with Less:} Developing ultra-high-performance concrete formulations to optimize structural geometry and reduce material volumes while achieving identical or superior load-bearing capacity.
    \item \textbf{Making Buildings Sustainable:} Implementing advanced integrated systems (such as specialized insulation, advanced roofing matrices, and high-durability thermal envelopes) to mitigate lifetime operational carbon.
    \item \textbf{Circular Construction (Building New from Old):} Developing local waste sorting loops to fully reuse construction demolition waste as raw kiln feedstock.
\end{enumerate}

\subsubsection{The "5C" Framework for Carbon Mitigation}
Historically, in 1990, the standard baseline carbon intensity stood at \num{783}~\si{\kilogram} of $\text{CO}_2$ per tonne of cement produced. To systematically eliminate this inventory, Holcim aligns its mitigation levers within the industry-recognized \textbf{5C Framework}:
\begin{itemize}
    \item \textbf{Clinker:} Lowering the traditional clinker-to-cement ratio (clinker factor) by introducing alternative supplementary cementitious components (SCMs) like slag or calcined clays. Holcim has successfully demonstrated a \num{40}\% clinker reduction across specific commercial lines.
    \item \textbf{Cement:} Developing low-emission blended formulations that achieve high strength with minimal embodied binder volumes.
    \item \textbf{Concrete:} Optimizing ready-mix concrete mix design matrices (e.g., the \MakeLowercase{\texttt{ECOPact}} product line, which yields \num{30}\% to \num{70}\% lower carbon footprints compared to ordinary Portland cement baselines).
    \item \textbf{Construction:} Promoting structural optimization via smart design, building information modeling (BIM), and modular construction to lower job-site material waste.
    \item \textbf{Carbon Capture and Sequestration (CCUS):} Capturing remaining process emissions directly from the flue gas stream using cutting-edge carbon purification systems.
\end{itemize}

% ==========================================
\subsection{Operational Optimization and Circular Asset Upgrades}
% ==========================================

\subsubsection{The Efficiency Shift: Horizontal Rotary Kilns to Vertical Preheater Towers}
A critical engineering milestone observed on-site is the decommissioning of traditional horizontal rotary dry/wet kilns in favor of modern, high-efficiency vertical preheater structures. 
\begin{itemize}
    \item \textbf{Geometric Architecture:} Horizontal rotary kilns require a continuous mechanical rotation of an inclined cylinder to move raw material against a counter-current hot air stream. This process is limited by high radiant heat losses and a lower specific surface area for heat exchange.
    \item \textbf{Vertical Tower Processing:} The raw material is introduced at the top of a multi-stage vertical suspension preheater tower (reaching up to \num{140}~\si{\meter} in height). It cascades downward through a series of cyclones, interacting directly with rising thermal exhaust gases. Gravity drives the material downward while it undergoes rapid calcination before entering a shortened rotary kiln for final sintering.
\end{itemize}

This structural transformation has optimized the Specific Thermal Energy Consumption (STEC). The transition \textbf{reduces energy requirements from \num{300}~\si{\mega\joule} down to \num{200}~\si{\mega\joule} per volumetric unit}, achieving a clear \num{40}\% net energy savings. Furthermore, modern dry processing eliminates the need for high water contents in the raw slurry. This removes the massive latent heat penalties associated with water evaporation, substantially improving global thermal efficiency.

\subsubsection{Alternative Fuels and the Role of Co-Processing}
Through its waste management subsidiary, \textit{Geocycle}, Holcim applies the principle of industrial \textbf{co-processing}, wherein waste streams simultaneously undergo material recycling and energy recovery inside the cement kiln. This operational strategy satisfies strict waste management hierarchies while targeting the total elimination of primary fossil fuels (reaching a thermal substitution rate $\text{TSR} > 95\%$ by 2030).

\begin{table}[!ht]
\centering
\small
\begin{tabular}{lp{11cm}}
\toprule
\textbf{Waste Stream Category} & \textbf{Physical Constitution and Engineering Function} \\
\midrule
\textbf{Fluff / SRF} & Solid Refined Fuels composed of shredded, non-recyclable industrial and municipal plastic fractions. Imparts high calorific content. \\
\textbf{Biomass} & Dried sewage sludges, animal meal, and impregnated sawdust. Provides carbon-neutral thermal energy. \\
\textbf{Liquid Hazardous Waste} & Spent industrial solvents, contaminated chemical emulsions, and waste oils. Handled via high-precision atomization lances. \\
\textbf{Pasty Waste} & Viscous paint sludges and chemical bottom residues. Destroyed safely under continuous high-temperature flame exposure. \\
\textbf{Mineral Waste} & Decarbonated calcium sources, steel slags, processing ashes, and contaminated soils. Handled via co-processing to supply essential $\text{Al}_2\text{O}_3$, $\text{Fe}_2\text{O}_3$, and $\text{SiO}_2$ ions directly to the clinker crystal phase. \\
\bottomrule
\end{tabular}
\caption{Geocycle waste stream co-processing inputs at the Belgian kilns.}
\label{tab:geocycle}
\end{table}

The total annual mass throughput handled by Geocycle inside Belgium includes \num{225}~\si{\kilo\tonne} (kt) of alternative fuels (including \num{36}~kt of liquid solvents), \num{130}~kt of alternative raw materials for clinker production, and \num{550}~kt of alternative mineral components for direct cement grinding. Crucially, the ultra-high operating temperatures ($\sim1400^\circ\text{C}$) and long residence times inside cement kilns provide complete thermal destruction of complex organic pollutants. This enables safe, clean material loops that can securely destroy hazardous substances such as per- and polyfluoroalkyl substances (PFAS).

\subsubsection{Ancillary Green Infrastructure}
To mitigate its Scope 2 footprint and decarbonize its logistics, Holcim Belgium has developed several supporting infrastructure projects:
\begin{itemize}
    \item \textbf{Floating Solar Array:} A large-scale solar photovoltaic plant constructed directly on an adjacent quarry lake to generate clean electricity.
    \item \textbf{Fleet Electrification:} Introduction of heavy-duty electric trucks and the deployment of Belgium’s first 100\% fully electric ready-mix concrete mixer truck.
\end{itemize}

% ==========================================
\subsection{Next-Generation Innovations: The GO4ZERO Project}
% ==========================================
The \MakeLowercase{\texttt{GO4ZERO}} initiative represents Holcim Belgium's flagship project to engineer a commercial net-zero cement plant at the Obourg site. Backed by a capital expenditure exceeding \num{550} million Euros, this multi-phase industrial transformation is scheduled across a rigid execution timeline:

\begin{table}[!ht]
\centering
\small
\begin{tabular}{llp{9cm}}
\toprule
\textbf{Phase} & \textbf{Target Date} & \textbf{Engineering Milestone and Operational Deliverable} \\
\midrule
\textbf{Phase I (a)} & Mid-2026 & Deployment of a new limestone crushing unit at Vaulx, extension of local mining permissions, and activation of low-emission rail transport loops. \\
\textbf{Phase I (b)} & January 2027 & Commissioning of an optimized \num{4500}~tonne-per-day (tpd) dry-process kiln operating in standard Air Mode for system stabilization and initial ramp-up. \\
\textbf{Oxyfuel Transition} & Mid-2027 & Verification trials of the air-oxyfuel switchable system and introductory oxy-combustion runs. \\
\textbf{Phase II} & Early 2029 & Full activation of the Cryogenic Purification Unit (CPU), permanent switch to Oxyfuel operation, and continuous carbon capture execution. \\
\textbf{Net-Zero Scale} & 2030 & Steady-state production of \num{2} million tonnes of fully decarbonized cement per year. \\
\bottomrule
\end{tabular}
\caption{GO4ZERO project industrial execution timeline.}
\label{tab:go4zero}
\end{table}

\subsubsection{Oxy-Combustion and Cryogenic Purification (CPU) Engineering}
Standard combustion using ambient air yields a flue gas stream highly diluted with nitrogen ($\sim 70-80\%\ \text{N}_2$), making carbon capture energetically inefficient. The \MakeLowercase{\texttt{GO4ZERO}} kiln re-engineers this process using \textbf{Oxyfuel technology}:
\begin{enumerate}
    \item \textbf{Oxygen Supply Line:} An industrial oxygen ($\text{O}_2$) pipeline intersects the center of the plant, delivering pure oxygen to the kiln burner.
    \item \textbf{Oxy-Combustion Mechanics:} Fossil and alternative fuels are combusted inside an atmosphere of pure $\text{O}_2$ blended with recycled carbon dioxide exhaust gas. This eliminates nitrogen from the combustion chamber.
    \item \textbf{Flue Gas Composition:} The resulting flue gas is highly concentrated, consisting almost entirely of carbon dioxide and water vapor ($\text{CO}_2 + \text{H}_2\text{O}$).
    \item \textbf{Cryogenic Purification (CPU):} The concentrated gas stream passes into an on-site CPU. The water is condensed out, and the gas is cooled to cryogenic temperatures under high pressure to liquefy the $\text{CO}_2$. This process isolates a pure, high-density liquid phase ready for transport and sequestration.
\end{enumerate}
This configuration will capture over \textbf{\num{1.1} million tonnes of $\text{CO}_2$ annually} starting in 2029, dropping the gross carbon footprint per tonne of cement from \num{566}~\si{\kilogram} down to a net-negative allocation after factoring in passive recarbonation.

\subsubsection{Raw Material Logistical Storage}
To manage feedstock variations without disrupting continuous kiln operations, a newly engineered storage hub has been built on-site. This includes a massive automated storage hangar and a specialized structural dome dedicated to raw limestone. This facility allows the simultaneous management of 5 distinct mineral materials, including 2 separate types of alternative pre-decarbonated materials. This ensures uniform chemical feedstocks even when incorporating high volumes of recycled components.

% ==========================================
\subsection{The Future of Materials: Circular Economy and Challenges}
% ==========================================

\subsubsection{Limits of Concrete Recycling}
While recycling construction waste is a key objective for the circular economy, complete closed-loop recycling faces major thermodynamic and chemical engineering challenges:
\begin{itemize}
    \item \textbf{Mechanical Separation:} Fine-grain concrete aggregates are highly adhered to old hydrated cement paste. Separating the original sand and gravel from this paste requires high mechanical energy.
    \item \textbf{The Micro-Mineral Challenge:} Reclaiming the reacted cement matrix to reuse it as a direct substitute for unreacted cement clinker is technically difficult. Hydrated cement consists of stable calcium-silicate-hydrates ($\text{C-S-H}$ gel) and calcium hydroxide ($\text{Ca(OH)}_2$). These phases have lost their hydraulic reactivity. To reactivate them, they must be re-calcined at high temperatures inside a kiln.
\end{itemize}
Consequently, the vast majority of current structural concrete recycling is limited to downcycling. This involves crushing old concrete to use as coarse road base or stabilization aggregate, rather than recovering pure cementitious material. To break through this limit, advanced sorting projects like \textit{Recygénie} in Paris are testing high-purity mineral sorting to develop concrete mixes where all components are derived from recycled materials.

\subsubsection{Regulatory and Cost Barriers}
The widespread industrial adoption of recycled binders faces significant hurdles:
\begin{enumerate}
    \item \textbf{Regulatory Latency:} Current construction structural codes (such as EN 197-1 and national standards) enforce strict limits on the mineralogical components allowed in structural concrete. While alternative low-carbon and recycled raw materials perform well in laboratory durability and mechanical tests, standard regulations have not yet evolved to permit wide-scale structural deployment.
    \item \textbf{Macro-Economic Cost Structures:} Although alternative and recycled materials offer clear environmental benefits, the logistics of sorting, processing, and purifying waste mean they remain more expensive than extracting primary natural raw materials.
\end{enumerate}

\subsubsection{Global Structural Demand Discrepancies}
The future consumption curve for cement follows divergent paths depending on regional economic development:
\begin{itemize}
    \item \textbf{Developing Nations:} Exhibiting high growth. Rapid urbanization, demographic expansion, and major infrastructure deficits ensure that demand for primary concrete will continue to expand across these regions.
    \item \textbf{Developed Nations (Europe / North America):} Exhibiting a plateau or structural decline. Population levels are stable, and primary infrastructure networks are largely complete. As a result, demand for volume is stabilizing. Over time, consumption is expected to contract slightly as alternative building materials, strict building norms, and circular design alternatives replace traditional primary binders.
\end{itemize}

\subsubsection{Projected Structural Cost Sensitivity}
A key consideration for project managers and structural designers is the cost elasticity of cement within the wider real estate market. Quantitative financial modeling indicates that \textbf{if the base market price of cement doubles, the total final construction cost of a building asset increases by only \num{5}\% to \num{10}\%}. Because raw material costs make up a relatively small share of total real estate development costs, low-carbon binders remain commercially viable. This minimal sensitivity means the industry can absorb the higher costs of deploying advanced carbon capture technologies without undermining the financial viability of downstream construction projects.